\documentclass[11pt,a4paper]{article}
\usepackage[utf8]{inputenc}
\usepackage[spanish]{babel}
\usepackage{amsmath, amssymb}
\usepackage[margin=2.5cm]{geometry}
\usepackage{fancyhdr}
\usepackage{tabularx}
\usepackage{booktabs}
\usepackage{tcolorbox}

\definecolor{UCBblue}{RGB}{0, 51, 153}
\definecolor{UCBgray}{RGB}{100, 100, 100}

\pagestyle{fancy}
\fancyhf{}
\fancyhead[L]{\textcolor{UCBblue}{\textbf{Ingeniería Mecatrónica - UCB Tarija}}}
\fancyhead[R]{\textcolor{UCBgray}{Guión de Clase Ajustado - Sesión 2}}
\fancyfoot[C]{\thepage}
\renewcommand{\headrulewidth}{0.4pt}
\renewcommand{\headrule}{\hbox to\headwidth{\color{UCBblue}\leaders\hrule height \headrulewidth\hfill}}

\begin{document}

\begin{center}
    {\Large \textbf{\textcolor{UCBblue}{PLANIFICACIÓN DIDÁCTICA Y GUIÓN DE CLASE}}}\\
    \vspace{0.3cm}
    {\large \textbf{Sesión 2: Consolidación Matricial e Introducción al Plano Complejo (CA)}}\\
    \vspace{0.5cm}
    \textbf{Asignatura:} Taller de Nivelación - Álgebra Lineal Aplicada a la Teoría de Redes \\
    \textbf{Docente:} M.Sc. Ing. Bernardo Quiroga Turdera \quad \textbf{Duración:} 120 Minutos
\end{center}

\vspace{0.3cm}

\section*{1. Estrategia de Ajuste Pedagógico (Transición Suave)}
Debido a la necesidad de nivelación algebraica basal ocurrida en la Sesión 1, esta segunda sesión se divide en dos fases técnicas de distinta naturaleza para proteger la carga cognitiva:
\begin{itemize}
    \item \textbf{Fase A (50 min):} Cierre de competencias computacionales en Corriente Continua (Ejecución de los Notebooks 01 y 02 pospuestos).
    \item \textbf{Fase B (70 min):} Apertura del Módulo 3. Fundamentación epistemológica de la necesidad de los números complejos en ingeniería de CA (\textit{Just-In-Time Physics}) y su sintaxis matricial primaria en Python.
\end{itemize}

\section*{2. Cronograma de Ejecución (120 Minutos)}
\begin{table}[h!]
    \centering
    \renewcommand{\arraystretch}{1.3}
    \begin{tabularx}{\textwidth}{@{} c c X l @{}}
        \toprule
        \textbf{Hora} & \textbf{Fase CDIO} & \textbf{Actividad Principal} & \textbf{Material de Soporte} \\
        \midrule
        17:00 - 17:25 & \textbf{Implementar} & \textit{Catch-up}: Co-Programación del tensor bidimensional de la red de potencia. Validación de simetría estructural. & Notebook 01 \\
        17:25 - 17:50 & \textbf{Operar} & Cómputo del determinante, fallas (\texttt{LinAlgError}) y Regla de Cramer asistida. & Notebook 02 \\
        \midrule
        17:50 - 18:20 & \textbf{Concebir} & Limitaciones de la CC. Componentes dinámicos, el desfase y el Plano Complejo ($s$). & Beamer (Intro CA) \\
        18:20 - 18:50 & \textbf{Diseñar} & El Fasor: Conversión trigonométrica analítica a pulso (Rectangular $\leftrightarrow$ Polar). & Guía Analítica (S2) \\
        18:50 - 19:00 & \textbf{Implementar} & Declaración de la constante nativa \texttt{1j} en Python. Extracción de magnitud y fase. & Notebook 03 \\
        \bottomrule
    \end{tabularx}
\end{table}

\section*{3. Consideraciones y Notas para el Docente}
\begin{itemize}
    \item \textbf{Gestión de la Frustración:} Durante la Fase A, enfatice que el código de Python es simplemente la "calculadora automática" del esfuerzo cognitivo que ellos ya realizaron a mano ayer en la pizarra.
    \item \textbf{El Concepto de "$j$":} En la Fase B (17:50), deténgase a explicar que en matemáticas abstractas el número imaginario se denomina $i$, pero en ingeniería eléctrica adoptamos la convención $j$ para evitar un choque de variables con la corriente eléctrica en el dominio del tiempo $i(t)$.
    \item \textbf{Cierre de la Sesión:} Al finalizar el Notebook 03, muestre a los estudiantes cómo Python maneja la división compleja internamente, demostrando que el álgebra vectorial ya no es un obstáculo mecánico, preparándolos para resolver matrices enteras con números complejos en la Sesión 3.
\end{itemize}

\end{document}
